% F7 -- Accuracy is dictated by how often the user changes cell.
\begin{figure}[t]
\centering
\replegend{\swatch{clbase}~persistence baseline \qquad \swatch{clmodel}~model}
\fitwidth{%
\begin{tikzpicture}
\begin{groupplot}[
  group style={group size=2 by 1, horizontal sep=1.5cm, ylabels at=edge left},
  repbar, width=0.47\linewidth, height=5.2cm,
  ymin=0, ymax=125, ytick={0,20,...,100},
  ylabel={top-1 accuracy (\%)},
  /pgf/bar width=7pt, xtick=data,
  x tick label style={font=\scriptsize,rotate=35,anchor=east},
  title style={font=\footnotesize,yshift=1pt},
]
\nextgroupplot[title={(a) \srawbig, five equal groups},
  symbolic x coords={qa,qb,qc,qd,qe},
  xticklabels={0\%,0--21\%,21--38\%,38--60\%,60--100\%},
  xlabel={share of steps where the cell changes}]
\addplot[fill=clbase,draw=none] coordinates {(qa,100.0) (qb,88.4) (qc,71.7) (qd,52.8) (qe,26.6)};
\addplot[fill=clmodel,draw=none] coordinates {(qa,100.0) (qb,87.6) (qc,71.6) (qd,56.5) (qe,34.4)};
\node[font=\scriptsize,text=clgain,anchor=west,rotate=90,yshift=7pt] at (axis cs:qd,74) {$+3.7\pt$};
\node[font=\scriptsize,text=clgain,anchor=west,rotate=90,yshift=7pt] at (axis cs:qe,52) {$+7.8\pt$};
%
\nextgroupplot[title={(b) \smob, four equal groups},
  symbolic x coords={pa,pb,pc,pd},
  xticklabels={0--25\%,25--50\%,50--75\%,75--100\%},
  xlabel={quartile of movement rate}]
\addplot[fill=clbase,draw=none] coordinates {(pa,87.7) (pb,79.6) (pc,66.7) (pd,42.8)};
\addplot[fill=clmodel,draw=none] coordinates {(pa,80.6) (pb,74.3) (pc,63.3) (pd,43.4)};
\node[font=\scriptsize,text=clloss,anchor=west,rotate=90,yshift=7pt] at (axis cs:pa,105) {$-7.1\pt$};
\node[font=\scriptsize,text=clloss,anchor=west,rotate=90,yshift=7pt] at (axis cs:pb,97) {$-5.3\pt$};
\node[font=\scriptsize,text=clloss,anchor=west,rotate=90,yshift=7pt] at (axis cs:pc,84) {$-3.4\pt$};
\node[font=\scriptsize,text=clgain,anchor=west,rotate=90,yshift=7pt] at (axis cs:pd,61) {$+0.6\pt$};
\end{groupplot}
\end{tikzpicture}}
\caption{Accuracy by cell-switch rate}
\label{fig:switch-rate}
\figtext{%
Accuracy is set by how repetitive the user already is, and the model's
contribution runs the other way. In both panels the users are sorted by their
cell-switch rate, the share of consecutive steps at which the cell changes,
and cut into equal groups, five on the left and four on the right, so
volatility increases towards the right; each group carries the persistence
baseline in grey-blue and the model in dark blue, scored on identical
positions, with the gap annotated above. (a)~On the unfiltered day the fifth
of users who never change cell are predicted at 100\,\% by both, because for
them the task is to notice that nothing is happening, and the most volatile
fifth plateaus at 34.4\,\%. That is a 65-point spread produced entirely by who
the users are, on a single trained model, with no hyperparameter involved. The
model overtakes the baseline only in the two hardest groups, by $+3.7$ and
$+7.8\pt$, and those two hold 18.5\,\% of the predictions between them, the
most volatile one alone holding 5.6\,\%: eight points on one twentieth of the
volume moves the aggregate by less than half a point. (b)~The same ordering on
the 100 held-out mobile users of \smob{}, and less comfortably. The model
\emph{loses} to the baseline in the three easier quartiles ($-7.1$, $-5.3$,
$-3.4\pt$) and only draws level in the most volatile one ($+0.6\pt$), which is
precisely how a model can be genuinely useful on hard users and still be worse
than a one-line rule overall. Baselines are recomputed per user, on the same
positions the model is scored on.}
\end{figure}
